% Generated by tools/make_body.py; edit generator, not this file.
% UNPROOFREAD OCR TRANSCRIPTION: page-image verification remains required.
\chapter{Conformal Mapping. Dirichlet's Problem}\label{chap:6}

% Original PDF physical page 245
\sourcepage{245}
In the geometrically oriented part of the theory of analytic functions the problem of conformal mapping\index{Conformal mapping|(}\index{Mapping!conformal|(}\index{Conformal mapping}\index{Mapping!conformal} plays a dominating role. Existence and uniqueness theorems permit us to define important analytic functions without resorting to analytic expressions, and geometric properties of the regions that are being mapped lead to analytic properties of the mapping function.

The Riemann mapping theorem\index{Riemann mapping theorem|(}\index{Riemann mapping theorem} deals with the mapping of one simply connected region onto another. We shall give a proof that leans on the theory of normal families. To handle the more difficult case of multiply connected regions we shall have to solve the Dirichlet problem, which is the boundary-value problem for the Laplace equation.

\setcounter{section}{0}\section{The Riemann Mapping Theorem}\label{sec:6.1}
We shall prove that the unit disk can be mapped conformally onto any simply connected region in the plane, other than the plane itself. This will imply that any two such regions can be mapped conformally onto each other, for we can use the unit disk as an intermediary step. The theorem is applied to polygonal regions, and in this case an explicit form for the mapping function is derived.

\setcounter{section}{1}\setcounter{subsection}{0}\subsection{Statement and Proof}\label{sec:6.1.1}
Although the mapping theorem was formulated by Riemann, its first successful proof was due to


% Original PDF physical page 246
\sourcepage{246}
P.\index{Mapping!schlicht|(}\index{Mapping!univalent|(}\index{Schlicht function|(} Koebe\index{Koebe P@Koebe, P.}.\hyperlink{fn:246:dagger}{\textsuperscript{\ensuremath{\dagger}}} The proof we shall present is a shorter variant of the original proof.

\begin{originalstatement}{Theorem}{1}\label{thm:6.1}
Given any simply connected region $\Omega$ which is not the whole plane, and a point $z_0 \in \Omega$, there exists a unique analytic function $f(z)$ in $\Omega$, normalized by the conditions $f(z_0) = 0$, $f'(z_0) > 0$, such that $f(z)$ defines a one-to-one mapping of $\Omega$ onto the disk $|w| < 1$.
\end{originalstatement}
The uniqueness is easily proved, for if $f_1$ and $f_2$ are two such functions, then $f_1[f_2^{-1}(w)]$ defines a one-to-one mapping of $|w| < 1$ onto itself. We know that such a mapping is given by a linear transformation $S$ (Chap.~\hyperref[chap:4]{4}, Sec.~\hyperref[sec:4.3.4]{3.4}, Ex. \hyperref[exercise:4.3.4.7.5]{5}). The conditions $S(0) = 0$, $S'(0) > 0$ imply $S(w) = w$; hence $f_1 = f_2$.

An analytic function $g(z)$ in $\Omega$ is said to be univalent if $g(z_1) = g(z_2)$ only for $z_1 = z_2$, in other words, if the mapping by $g$ is one to one (the German word schlicht, which lacks an adequate translation, is also in common use). For the existence proof we consider the family $\mathfrak{F}$ formed by all functions $g$ with the following properties: (i) $g$ is analytic and univalent in $\Omega$, (ii) $|g(z)| \leq 1$ in $\Omega$, (iii) $g(z_0) = 0$ and $g'(z_0) > 0$. We contend that $f$ is the function in $\mathfrak{F}$ for which the derivative $f'(z_0)$ is a maximum. The proof will consist of three parts: (1) it is shown that the family $\mathfrak{F}$ is not empty; (2) there exists an $f$ with maximal derivative; (3) this $f$ has the desired properties.

To prove that $\mathfrak{F}$ is not empty we note that there exists, by assumption, a point $a \neq \infty$ not in $\Omega$. Since $\Omega$ is simply connected, it is possible to define a single-valued branch of $\sqrt{z - a}$ in $\Omega$; denote it by $h(z)$. This function does not take the same value twice, nor does it take opposite values. The image of $\Omega$ under the mapping $h$ covers a disk $|w - h(z_0)| < \rho$, and therefore it does not meet the disk $|w + h(z_0)| < \rho$. In other words, $|h(z) + h(z_0)| \geq \rho$ for $z \in \Omega$, and in particular $2|h(z_0)| \geq \rho$. It can now be verified that the function
\begin{equation*}
g _ {0} (z) = \frac {\rho}{4} \frac {| h ^ {\prime} (z _ {0}) |}{| h (z _ {0}) | ^ {2}} \cdot \frac {h (z _ {0})}{h ^ {\prime} (z _ {0})} \cdot \frac {h (z) - h (z _ {0})}{h (z) + h (z _ {0})}
\end{equation*}
belongs to the family $\mathfrak{F}$. Indeed, because it is obtained from the univalent function\index{Univalent function} h by means of a linear transformation, it is itself univalent. Moreover, $g_{0}(z_{0}) = 0$ and $g_{0}^{\prime}(z_{0}) = (\rho/8)|h^{\prime}(z_{0})|/|h(z_{0})|^{2} > 0$. Finally, the estimate
\begin{equation*}
\left| \frac {h (z) - h \left(z _ {0}\right)}{h (z) + h \left(z _ {0}\right)} \right| = | h \left(z _ {0}\right) | \cdot \left| \frac {1}{h \left(z _ {0}\right)} - \frac {2}{h (z) + h \left(z _ {0}\right)} \right| \leqslant \frac {4 \left| h \left(z _ {0}\right) \right|}{\rho}
\end{equation*}
shows that $|g_0(z)| \leq 1$ in $\Omega$.

\begingroup\renewcommand{\thefootnote}{}\footnotetext{\hypertarget{fn:246:dagger}{}\textsuperscript{\ensuremath{\dagger}} A related theorem from which the mapping theorem can be derived had been proved earlier by W. F. Osgood\index{Osgood W F@Osgood, W. F.|hyperpagenote}, but did not attract the attention it deserves.}\endgroup

% Original PDF physical page 247
\index{Mapping!schlicht|)}\index{Mapping!univalent|)}\index{Schlicht function|)}\sourcepage{247}
The derivatives $g'(z_0), g \in \mathfrak{F}$, have a least upper bound $B$ which a priori could be infinite. There is a sequence of functions $g_n \in \mathfrak{F}$ such that $g_n'(z_0) \to B$. By Chap.~\hyperref[chap:5]{5}, Theorem~\hyperref[thm:5.12]{12} the family $\mathfrak{F}$ is normal. Hence there exists a subsequence $\{g_{n_k}\}$ which tends to an analytic limit function $f$, uniformly on compact sets. It is clear that $|f(z)| \leq 1$ in $\Omega, f(z_0) = 0$ and $f'(z_0) = B$ (this proves that $B < +\infty$). If we can show that $f$ is univalent, it will follow that $f$ is in $\mathfrak{F}$ and has a maximal derivative at $z_0$.

In the first place $f$ is not a constant, for $f'(z_0) = B > 0$. Choose a point $z_1 \in \Omega$, and consider the functions $g_1(z) = g(z) - g(z_1)$, $g \in \mathfrak{F}$. They are all $\neq 0$ in the region obtained by omitting $z_1$ from $\Omega$. By Hurwitz's theorem (Chap.~\hyperref[chap:5]{5}, Theorem~\hyperref[thm:5.2]{2}) every limit function is either identically zero or never zero. But $f(z) - f(z_1)$ is a limit function, and it is not identically zero. Hence $f(z) \neq f(z_1)$ for $z \neq z_1$, and since $z_1$ was arbitrary we have proved that $f$ is univalent.

It remains to show that f takes every value w with $|w| < 1$. Suppose it were true that $f(z) \neq w_{0}$ for some $w_{0}, |w_{0}| < 1$. Then, since $\Omega$ is simply connected, it is possible to define a single-valued branch of
\begin{equation*}
F (z) = \sqrt {\frac {f (z) - w _ {0}}{1 - \bar {w} _ {0} f (z)}}\tag{1}\label{eq:6.1}
\end{equation*}
(Recall that all closed curves in a simply connected region are homologous to 0. If $\varphi(z) \neq 0$ in $\Omega$ we can define $\log \varphi(z)$ by integration of $\varphi'(z)/\varphi(z)$, and $\sqrt{\varphi(z)} = \exp\left(\frac{1}{2} \log \varphi(z)\right)$.)

It is clear that $F$ is univalent and that $|F| \leq 1$. To normalize it we form
\begin{equation*}
G (z) = \frac {\left| F ^ {\prime} (z _ {0}) \right|}{F ^ {\prime} (z _ {0})} \cdot \frac {F (z) - F (z _ {0})}{1 - \overline {{F (z _ {0})}} F (z)}\tag{2}\label{eq:6.2}
\end{equation*}
which vanishes and has a positive derivative at $z_{0}$. For its value we find, after brief computation,
\begin{equation*}
G ^ {\prime} (z _ {0}) = \frac {| F ^ {\prime} (z _ {0}) |}{1 - | F (z _ {0}) | ^ {2}} = \frac {1 + | w _ {0} |}{2 \sqrt {| w _ {0} |}} B > B.
\end{equation*}
This is a contradiction, and we conclude that $f(z)$ assumes all values $w$, $|w| < 1$. The proof is now complete.

At first glance it may seem like pure luck that our computation yields $G'(z_{0}) > f'(z_{0})$. This is not quite so, for the formulas \hyperref[eq:6.1]{(1)} and \hyperref[eq:6.2]{(2)} permit us to express $f(z)$ as a single-valued analytic function of $W = G(z)$ which maps $|W| < 1$ into itself. The inequality $|f'(z_{0})| < |G'(z_{0})|$ is therefore a consequence of Schwarz's lemma.

The purely topological content of Theorem~\hyperref[thm:6.1]{1} is important by itself. We know now that any simply connected region can be mapped topolog%
% Original PDF physical page 248
\sourcepage{248}%
ically onto a disk (for the whole plane a very simple mapping can be constructed), and hence any two simply connected regions are topologically equivalent.

\subsection*{EXERCISES}
\originalnumber{1}{exercise:6.1.1.1.1}
1. If $z_{0}$ is real and $\Omega$ is symmetric with respect to the real axis, prove by the uniqueness that $f$ satisfies the symmetry relation $f(\bar{z}) = \overline{f(z)}$.

\originalnumber{2}{exercise:6.1.1.1.2}
2. What is the corresponding conclusion if $\Omega$ is symmetric with respect to the point $z_0$?

\setcounter{section}{1}\setcounter{subsection}{1}\subsection{Boundary Behavior}\label{sec:6.1.2}
We are assuming that $f(z)$ defines a conformal mapping\index{Conformal mapping}\index{Mapping!conformal} of a region $\Omega$ onto another region $\Omega'$. What happens when $z$ approaches the boundary? There are cases where the boundary behavior\index{Boundary!behavior} can be foretold with great precision. For instance, if $\Omega$ and $\Omega'$ are Jordan regions,\hyperlink{fn:248:dagger}{\textsuperscript{\ensuremath{\dagger}}} then $f$ can be extended to a topological mapping of the closure of $\Omega$ onto the closure of $\Omega'$. Unfortunately, considerations of space do not permit us to include a proof of this important theorem (the proof would require a considerable amount of preparation).

What we can and shall prove is a very modest theorem of purely topological content. Let us first make it clear what we mean when we say that $z$ approaches the boundary of $\Omega$. There are two cases: we may consider a sequence $\{z_n\}$ of points in $\Omega$, or we may consider an arc $z(t)$, $0 \leq t < 1$, such that all $z(t)$ are in $\Omega$. We shall say that the sequence or the arc tends to the boundary if the points $z_n$ or $z(t)$ will ultimately stay away from any point in $\Omega$. In other words, if $z \in \Omega$ there shall exist an $\varepsilon > 0$ and an $n_0$ or a $t_0$ such that $|z_n - z| \geq \varepsilon$ for $n > n_0$, or such that $|z(t) - z| \geq \varepsilon$ for all $t > t_0$.

In this situation, the disks of center z and radius $\varepsilon$ (which may depend on z) form an open covering of $\Omega$. It follows that any compact subset $K \subset \Omega$ is covered by a finite number of these disks. If we consider the largest of the corresponding $n_{0}$ or $t_{0}$ we find that $z_{n}$ or $z(t)$ cannot belong to K for $n > n_{0}$ or $t > t_{0}$. Colloquially speaking, for any compact set $K \subset \Omega$ there exists a tail end of the sequence or of the arc which does not meet K. Conversely, this implies the original condition, for if $z \in \Omega$ is given we may choose K to be a closed disk with center z that is contained in $\Omega$. If the radius of the disk is $\rho$ the original statement holds for any $\varepsilon < \rho$.

After these preliminary considerations the theorem we shall prove is almost trivial:

\begingroup\renewcommand{\thefootnote}{}\footnotetext{\hypertarget{fn:248:dagger}{}\textsuperscript{\ensuremath{\dagger}} It is known, although not so easy to prove, that a Jordan curve (Chap.~\hyperref[chap:3]{3}, Sec.~\hyperref[sec:3.2.1]{2.1}) divides the plane into exactly two regions, one bounded and one unbounded. The bounded region is called a Jordan region.}\endgroup

% Original PDF physical page 249
\sourcepage{249}
\begin{originalstatement}{Theorem}{2}\label{thm:6.2}
Let f be a topological mapping of a region $\Omega$ onto a region $\Omega'$. If $\{z_{n}\}$ or $z(t)$ tends to the boundary of $\Omega$, then $\{f(z_{n})\}$ or $f(z(t))$ tends to the boundary of $\Omega'$.
\end{originalstatement}
Indeed, let $K$ be a compact set in $\Omega'$. Then $f^{-1}(K)$ is a compact set in $\Omega$, and there exists $n_0$ (or $t_0$) such that $z_n$ (or $z(t)$) is not in $f^{-1}(K)$ for $n > n_0$ (or $t > t_0$). But then $f(z_n)$ [or $f(z(t))$] is not in $K$.

Although the theorem is topological, it is the application to conformal mappings\index{Conformal mapping}\index{Mapping!conformal} that is of greatest interest to us.

\setcounter{section}{1}\setcounter{subsection}{2}\subsection{Use of the Reflection Principle}\label{sec:6.1.3}
Stronger statements become possible if we have more information. We are mainly interested in simply connected regions and may therefore assume that one of the regions is a disk. With the same notation as in Sec.~\hyperref[sec:6.1.1]{1.1}, let $f(z)$ define a conformal mapping of the region $\Omega$ onto the unit disk with the normalization $f(z_{0}) = 0$ (the normalization by the derivative is irrelevant). We shall derive additional information by use of the reflection principle (Chap.~\hyperref[chap:4]{4}, Theorem 26).

Let us assume that the boundary of $\Omega$ contains a segment $\gamma$ of a straight line. Because rotations are unimportant we may as well suppose that $\gamma$ lies on the real axis; let it be the interval a < x < b. The assumption involves a significant simplification only if the rest of the boundary stays away from $\gamma$. For this reason we shall strengthen the hypothesis and require that every point of $\gamma$ has a neighborhood whose intersection with the whole boundary $\partial\Omega$ is the same as its intersection with $\gamma$. We say then that $\gamma$ is a free boundary arc.

By this assumption every point on $\gamma$ is the center of a disk whose intersection with $\partial\Omega$ is its real diameter. It is clear that each of the half disks determined by this diameter is entirely in or entirely outside of $\Omega$, and at least one must be inside. If only one is inside we call the point a one-sided boundary point, and if both are inside it is a two-sided boundary point. Because $\gamma$ is connected all its points will be of the same kind, and we speak of a one-sided or a two-sided boundary arc.

\begin{originalstatement}{Theorem}{3}\label{thm:6.3}
Suppose that the boundary of a simply connected region $\Omega$ contains a line segment $\gamma$ as a one-sided free boundary arc. Then the function $f(z)$ which maps $\Omega$ onto the unit disk can be extended to a function which is analytic and one to one on $\Omega \cup \gamma$. The image of $\gamma$ is an arc $\gamma'$ on the unit circle.
\end{originalstatement}
For two-sided arcs the same will be true with obvious modifications.

For the proof we consider a disk around $x_0 \in \gamma$ which is so small that the half disk in $\Omega$ does not contain the point $z_0$ with $f(z_0) = 0$. Then
% Original PDF physical page 250
\sourcepage{250}
log $f(z)$ has a single-valued branch in the half disk, and its real part tends to 0 as z approaches the diameter, for we know by Theorem~\hyperref[thm:6.2]{2} that $|f(z)|$ tends to 1. It follows by the reflection principle that $\log f(z)$ has an analytic extension to the whole disk. Therefore $\log f(z)$, and consequently $f(z)$, is analytic at $x_{0}$. The extensions to overlapping disks must coincide and define a function which is analytic on $\Omega \cup \gamma$.

We note further that $f'(z) \neq 0$ on $\gamma$. Indeed, $f'(x_0) = 0$ would imply that $f(x_0)$ were a multiple value, in which case the two subarcs of $\gamma$ that meet at $x_0$ would be mapped on arcs that form an angle $\pi / n$ with $n \geq 2$; this is clearly impossible. If, for instance, the upper half disks are in $\Omega$, then
\begin{equation*}
\partial \log | f | / \partial y = - \partial \arg f / \partial x <   0
\end{equation*}
on $\gamma$, and $\arg f$ moves constantly in the same direction. This proves that the mapping is one to one on $\gamma$.

The theorem can be generalized to regions with free boundary arcs on a circle. With obvious modifications the theorem is also true for two-sided boundary arcs.

\setcounter{section}{1}\setcounter{subsection}{3}\subsection{Analytic Arcs}\label{sec:6.1.4}
A real or complex function $\varphi(t)$ of a real variable $t$, defined on an interval $a < t < b$, is said to be real analytic (or analytic in the real sense) if, for every $t_0$ in the interval, the Taylor development $\varphi(t) = \varphi(t_0) + \varphi'(t_0)(t - t_0) + \frac{1}{2}\varphi'(t_0)(t - t_0)^2 + \cdots$ converges in some interval $(t_0 - \rho, t_0 + \rho)$, $\rho > 0$. But if this is so we know by Abel's theorem that the series is also convergent for complex values of $t$, as long as $|t - t_0| < \rho$, and that it represents an analytic function in that disk. In overlapping disks the functions are the same, for they coincide on a segment of the real axis. We conclude that $\varphi(t)$ can be defined as an analytic function in a region $\Delta$, symmetric to the real axis, which contains the segment $(a,b)$.

In these circumstances we say that $\varphi(t)$ determines an analytic arc\index{Analytic arc}\index{Arc!analytic}. It is regular if $\varphi'(t) \neq 0$, and it is a simple arc if $\varphi(t_1) = \varphi(t_2)$ only when $t_1 = t_2$.

We shall assume that the boundary of $\Omega$ contains a regular, simple, analytic arc $\gamma$, and that it is a free one-sided arc. The definition could be modeled on the previous one, but to avoid long explanations we shall assume offhand that there exists a region $\Delta$, symmetric to the interval $(a,b)$, with the property that $\varphi(t) \in \Omega$ when t lies in the upper half of $\Delta$, and that $\varphi(t)$ lies outside $\Omega$ for t in the lower half.

If $f(z)$ is the mapping function with $f(z_0) = 0$, and if we take care that $\varphi(t) \neq z_0$ in $\Delta$, then the reflection principle tells us that $\log f(\varphi(t))$, and hence $f(\varphi(t))$, has an analytic extension from the upper to the lower half of $\Delta$. For a real $t_0 \in (a, b)$ we know further that $\varphi'(t_0) \neq 0$. There%
% Original PDF physical page 251
\sourcepage{251}%
fore\index{Polygon!conformal mapping of|(} $\varphi$ has an analytic inverse $\varphi^{-1}$ in a neighborhood of $\varphi(t_{0})$, and it follows by composition that $f(z)$ is analytic in that neighborhood.

\begin{originalstatement}{Theorem}{4}\label{thm:6.4}
If the boundary of $\Omega$ contains a free one-sided analytic arc $\gamma$, then the mapping function has an analytic extension to $\Omega \cup \gamma$, and $\gamma$ is mapped on an arc of the unit circle.
\end{originalstatement}
We trust the reader to make the last statement more precise and to complete the proof.

\setcounter{section}{1}\section{Conformal Mapping Of Polygons}\index{Conformal mapping}\index{Mapping!conformal}\label{sec:6.2}
When $\Omega$ is a polygon, the mapping problem has an almost explicit solution. Indeed, we shall find that the mapping function can be expressed through a formula in which only certain parameters have values that depend on the specific shape of the polygon.

\setcounter{section}{2}\setcounter{subsection}{0}\subsection{The Behavior at an Angle}\label{sec:6.2.1}
We assume that $\Omega$ is a bounded simply connected region whose boundary is a closed polygonal line without self-intersections. Let the consecutive vertices be $z_{1},\ldots ,z_{n}$ in positive cyclic order (we set $z_{n + 1} = z_1$). The angle at $z_{k}$ is given by the value of $\arg (z_{k - 1} - z_k) / (z_{k + 1} - z_k)$ between 0 and $2\pi$. We shall denote it by $\alpha_{k}\pi,0 < \alpha_{k} < 2$. It is also convenient to introduce the outer angles $\beta_{k}\pi = (1 - \alpha_{k})\pi, - 1 < \beta_{k} < 1$. Observe that $\beta_{1} + \dots +\beta_{n} = 2$. The polygon is convex if and only if all $\beta_{k} > 0$.

We know by Theorem~\hyperref[thm:6.3]{3} that the mapping function $f(z)$ can be extended by continuity to any side of the polygon (that is, to the open line segment between two consecutive vertices), and that each side is mapped in a one-to-one way onto an arc of the unit circle. We wish to show that these arcs are disjoint and leave no gap between them.

To see this we consider a circular sector $S_{k}$ which is the intersection of $\Omega$ with a sufficiently small disk about $z_{k}$. A single-valued branch of $\zeta = (z - z_k)^{1 / \alpha_k}$ maps $S_{k}$ onto a half disk $S_{k}'$. A suitable branch of $z_{k} + \zeta^{\alpha_{k}}$ has its values in $\Omega$, and we may consider the function $g(\zeta) = f(z_{k} + \zeta^{\alpha_{k}})$ in $S_{k}'$. It follows by Theorem~\hyperref[thm:6.2]{2} that $|g(\zeta)| \to 1$ as $\zeta$ approaches the diameter. The reflection principle applies, and we conclude that $g(\zeta)$ has an analytic continuation to the whole disk. In particular, this implies that $f(z)$ has a limit $w_{k} = e^{i\theta_{k}}$ as $z \to z_{k}$, and we find that the arcs that correspond to the sides meeting at $z_{k}$ do indeed have a common end point. Since $\arg f(z)$ must increase as $z$ traces the boundary in positive direction, the arcs do not overlap, at least not in a neighborhood of $w_{k}$. If we take into account that $f$ maps the boundary


% Original PDF physical page 252
\index{Riemann mapping theorem|)}\sourcepage{252}
\bookfigure{252}{1}{FIG. 6-1. Conformal mapping of a polygon.}
on a curve with winding number 1 about the origin, it can easily be concluded that all the arcs are mutually disjoint. In other words, f can be extended to a homeomorphic map of $\Omega^{-}$ onto the closed unit disk, the vertices $z_{k}$ go into points $w_{k}$, and the sides correspond to the arcs between these points (Fig.~\hyperref[fig:6-1]{6-1}).

\setcounter{section}{2}\setcounter{subsection}{1}\subsection{The Schwarz-Christoffel Formula}\index{Schwarz-Christoffel formula|(}\index{Schwarz-Christoffel formula}\label{sec:6.2.2}
The formula we are looking for refers not to the function f, but to its inverse function, which we shall denote by F.

\begin{originalstatement}{Theorem}{5}\label{thm:6.5}
The functions $z = F(w)$ which map $|w| < 1$ conformally onto polygons with angles $\alpha_k\pi\,(k = 1, \ldots, n)$ are of the form
\begin{equation*}
F (w) = C \int_ {0} ^ {w} \prod_ {k = 1} ^ {n} (w - w _ {k}) ^ {- \beta_ {k}} d w + C ^ {\prime}\tag{3}\label{eq:6.3}
\end{equation*}
where $\beta_{k} = 1 - \alpha_{k}$, the $w_{k}$ are points on the unit circle, and $C$, $C'$ are complex constants.

\end{originalstatement}
Because the function $g(\zeta) = f(z_k + \zeta^{\alpha_k})$ considered in the last paragraph of 2.1 is analytic at the origin, it has a Taylor development
\begin{equation*}
f (z _ {k} + \zeta^ {\alpha_ {k}}) = w _ {k} + \sum_ {m = 1} ^ {\infty} a _ {m} \zeta^ {m}.
\end{equation*}
Here $a_1 \neq 0$, for otherwise the image of the half disk $S_k'$ could not be contained in the unit disk. Therefore the series can be inverted, and on setting $w = f(z_k + \zeta^{\alpha_k})$ we obtain
\begin{equation*}
\zeta = \sum_ {m = 1} ^ {\infty} b _ {m} (w - w _ {k}) ^ {m}
\end{equation*}
% Original PDF physical page 253
\sourcepage{253}
with $b_{1} \neq 0$, the development being valid in a neighborhood of $w_{k}$. We raise to the power $\alpha_{k}$ and find, in terms of the inverse function F,
\begin{equation*}
F (w) - z _ {k} = (w - w _ {k}) ^ {\alpha_ {k}} G _ {k} (w)
\end{equation*}
where $G_{k}$ is analytic and $\neq0$ near $w_{k}$. It follows by differentiation that $F'(w)(w-w_{k})^{\beta_{k}}$ is analytic and $\neq0$ at $w_{k}$, and therefore the product
\begin{equation*}
H (w) = F ^ {\prime} (w) \prod_ {k = 1} ^ {n} (w - w _ {k}) ^ {\beta_ {k}}\tag{4}\label{eq:6.4}
\end{equation*}
is analytic and $\neq0$ in the closed unit disk.

We claim that $H(w)$ is, in actual fact, a constant. For this purpose we examine its argument when $w = e^{i\theta}$ lies on the unit circle between $w_k = e^{i\theta_k}$ and $w_{k+1} = e^{i\theta_{k+1}}$. We know that $\arg F'(e^{i\theta})$ equals the angle between the tangent to the unit circle at $e^{i\theta}$ and the tangent to its image at $F(e^{i\theta})$; with an abbreviated notation we express this by $\arg F' = \arg dF - \arg dw$. But $\arg dF$ is constant because $F$ describes a straight line, and $\arg dw = \theta + \pi/2$. The factor $w - w_k$ can be written $e^{i\theta} - e^{i\theta_k} = 2ie^{i(\theta + \theta_k)/2} \sin \frac{1}{2} (\theta - \theta_k)$, and hence its argument is $\theta/2$ plus a constant (this is also evident geometrically). When we add the arguments of all factors on the right-hand side of \hyperref[eq:6.4]{(4)} we find that $\arg H(w)$ differs by a constant from $-\theta + \left( \sum_{k=1}^{n} \beta_k \right) \cdot \theta/2 = 0$. Thus we conclude that $\arg H(w)$ is constant between $w_k$ and $w_{k+1}$, and since it is continuous it must be constant on the whole unit circle. The maximum principle permits us to conclude that $\arg H(w) = \operatorname{Im} \log H(w)$ is constant inside the unit circle, and so is consequently $H(w)$.

We have now proved that
\begin{equation*}
F ^ {\prime} (w) = C \prod_ {k = 1} ^ {n} (w - w _ {k}) ^ {- \beta_ {k}},
\end{equation*}
and formula \hyperref[eq:6.3]{(3)} follows by integration.

We remark that a linear transformation of the unit circle permits us to place three of the points $w_{k}$, for instance, $w_{1}$, $w_{2}$, $w_{3}$, in prescribed positions. For n = 3 we see that the mapping function depends only on the angles, except for trivial variable transformations; this reflects the fact that triangles with the same angles are similar. For n > 3 the remaining constants $w_{4}$, $\ldots$, $w_{n}$, or their arguments $\theta_{k}$, are called the accessory parameters\index{Accessory parameter} of the problem. It is only in rare cases that they can be determined other than by numerical computation.

If we give arbitrary values to the $\theta_{k}$ it is quite easy to verify that a function of the form \hyperref[eq:6.3]{(3)} maps the unit circle on a closed polygonal line,
% Original PDF physical page 254
\sourcepage{254}
but usually we are unable to tell whether it will intersect itself or not. If it does not, it is not difficult to show that $F(w)$, as given by \hyperref[eq:6.3]{(3)}, yields a one-to-one mapping onto the inside of the polygonal line (the precise proof makes use of the argument principle).

Formula \hyperref[eq:6.3]{(3)} is known as the Schwarz-Christoffel formula\index{Schwarz-Christoffel formula}. Another version of the same formula serves to map the upper half plane onto the inside of a polygon. The mapping function, from $\operatorname{Im} w > 0$ to $\Omega$, can now be written in the form
\begin{equation*}
F (w) = C \int_ {0} ^ {w} \prod_ {k = 1} ^ {n - 1} (w - \xi_ {k}) ^ {- \beta_ {k}} d w + C ^ {\prime}\tag{5}\label{eq:6.5}
\end{equation*}
where the $\xi_{k}$ are real. The last exponent $\beta_{n}$ does not appear explicitly in the formula, but it is determined by $\beta_{n}=2-(\beta_{1}+\cdots+\beta_{n-1})$, and like the other exponents it is subject to the condition $-1<\beta_{n}<1$. It then follows that the integral \hyperref[eq:6.5]{(5)} converges for $w=\infty$, and the point at $\infty$ will correspond to a vertex with angle $\alpha_{n}\pi$, $\alpha_{n}=1-\beta_{n}$. If $\beta_{n}=0$ the vertex is only apparent, and the polygon reduces to one with n-1 sides.

\subsection*{EXERCISES}
\originalnumber{1}{exercise:6.2.2.2.1}
1. Show that the $\beta_{k}$ in \hyperref[eq:6.3]{(3)} may be allowed to become $= -1$. What is the geometric interpretation?

\originalnumber{2}{exercise:6.2.2.2.2}
2. If a vertex of the polygon is allowed to be at $\infty$, what modification does the formula undergo? If in this context $\beta_{k}=1$, what is the polygon like?

\originalnumber{3}{exercise:6.2.2.2.3}
3. Show that the mappings of a disk onto a parallel strip, or onto a half strip with two right angles, can be obtained as special cases of the Schwarz-Christoffel formula.

\originalnumber{4}{exercise:6.2.2.2.4}
4. Derive formula \hyperref[eq:6.5]{(5)}, either directly or with the help of \hyperref[eq:6.3]{(3)}.

\originalnumber{5}{exercise:6.2.2.2.5}
5. Show that
\begin{equation*}
F (w) = \int_ {0} ^ {w} (1 - w ^ {n}) ^ {- 2 / n} d w
\end{equation*}
maps $|w| < 1$ onto the interior of a regular polygon with $n$ sides.

\originalnumber{6}{exercise:6.2.2.2.6}
6. Determine a conformal mapping\index{Conformal mapping}\index{Mapping!conformal} of the upper half plane on the region $\Omega = \{z = x + iy; x > 0, y > 0, \min(x, y) < 1\}$.

\setcounter{section}{2}\setcounter{subsection}{2}\subsection{Mapping on a Rectangle}\index{Rectangle, mapping on a|(}\index{Rectangle, mapping on a}\label{sec:6.2.3}
In case $\Omega$ is a rectangle we may choose $x_{1} = 0$, $x_{2} = 1$, $x_{3} = \rho > 1$ in \hyperref[eq:6.5]{(5)}. The mapping function will thus be given by
\begin{equation*}
F (w) = \int_ {0} ^ {w} \frac {d w}{\sqrt {w (w - 1) (w - \rho)}}
\end{equation*}
% Original PDF physical page 255
\index{Schwarz-Christoffel formula|)}\sourcepage{255}
which is an elliptic integral\index{Elliptic integral}. To be unambiguous we decide that the values of $\sqrt{w}$, $\sqrt{w-1}$, and $\sqrt{w-\rho}$ shall lie in the first quadrant. For a detailed study of the mapping, let us follow $F(w)$ as w traces the real axis. When w is real, each of the square roots is either positive or purely imaginary with a positive imaginary part (save for the point where the square root is 0). As 0 < w < 1 there are one real and two imaginary square roots. Therefore $F(w)$ decreases from 0 to a value -K where
\begin{equation*}
K = \int_ {0} ^ {1} \frac {d t}{\sqrt {t (1 - t) (\rho - t)}}.\tag{6}\label{eq:6.6}
\end{equation*}
For $1 < w < \rho$ there is only one imaginary square root. It follows that the integral from 1 to w is purely imaginary with a negative imaginary part. Hence $F(w)$ will follow a vertical segment from -K to -K - iK',
\begin{equation*}
K ^ {\prime} = \int_ {1} ^ {\rho} \frac {d t}{\sqrt {t (t - 1) (\rho - t)}}.
\end{equation*}
For $w > \rho$ the integrand is positive, and $F(w)$ will trace a horizontal segment in the positive direction. How far does it extend? Since the image is to be a rectangle, it must end at the point $-iK'$, but we prefer a direct verification. One way is to express the length of the segment by the integral
\begin{equation*}
\int_ {\rho} ^ {\infty} \frac {d t}{\sqrt {t (t - 1) (t - \rho)}}
\end{equation*}
and to show by the change of integration variable $t = (\rho - u)/(1 - u)$ that the integral transforms to \hyperref[eq:6.6]{(6)}. It is easier, however, to observe that Cauchy's theorem yields
\begin{equation*}
\int_ {- \infty} ^ {\infty} \frac {d t}{\sqrt {t (t - 1) (t - \rho)}} = 0,
\end{equation*}
for the integral over a semicircle with radius R tends to 0 as $R \to \infty$. The vanishing of the real part implies the equality of the horizontal segments, and from the vanishing of the imaginary part we deduce that $-\infty < w < 0$ is mapped on the segment from $-iK'$ to 0. The rectangle is completed.

It is often preferable to use a formula which reflects the double symmetry of the rectangle. The vertices can be made to correspond to points $\pm1$ and $\pm1/k$ with 0 < k < 1. Since a constant factor does not matter we can choose the mapping to be given by
\begin{equation*}
F (w) = \int_ {0} ^ {w} \frac {d w}{\sqrt {(1 - w ^ {2}) (1 - k ^ {2} w ^ {2})}},\tag{7}\label{eq:6.7}
\end{equation*}
% Original PDF physical page 256
\sourcepage{256}
\bookfigure{256}{1}{FIG. 6-2}
and this time we agree that $\sqrt{1-w^{2}}$ and $\sqrt{1-k^{2}w^{2}}$ shall have positive real parts. It is seen that the rectangle will have vertices at $-\frac{K}{2}$, $\frac{K}{2}$, $\frac{K}{2} + iK'$, $-\frac{K}{2} + iK'$ where
\begin{equation*}
\begin{array}{l} K = \int_ {- 1} ^ {1} \frac {d t}{\sqrt {(1 - t ^ {2}) (1 - k ^ {2} t ^ {2})}} \\ K ^ {\prime} = \int_ {1} ^ {1 / k} \frac {d t}{\sqrt {(t ^ {2} - 1) (1 - k ^ {2} t ^ {2})}}. \end{array}
\end{equation*}
The image of the upper half plane is the shaded rectangle $R_{0}$ in Fig.~\hyperref[fig:6-2]{6-2}. We denote the inverse function of F by $w = f(z)$; it is defined in $R_{0}$ and can be extended by continuity to a one-to-one mapping of the closed rectangle onto the closed half plane (with the topology of the Riemann sphere). Observe that $z = iK'$ corresponds to $\infty$.

The reflection principle allows us to extend the definition of $f$ to the adjacent rectangles $R_{1}$ and $R_{2}$, namely by setting $f(z) = \overline{f(\bar{z})}$ for $z \in R_{1}$ and $f(z) = \overline{f(K - \bar{z})}$ for $z \in R_{2}$. Similarly we can pass to $R_{3}$ either from $R_{1}$ or $R_{2}$; the extension is given by $f(z) = f(K - z)$. The process of reflection can obviously be continued until $f(z)$ is defined as a meromorphic function in the whole plane. It is perhaps even more convenient to define the extension by periodicity, for we find that the extended function must satisfy $f(z + 2K) = f(z)$, $f(z + 2iK') = f(z)$.

We have shown that the inverse function of the elliptic integral \hyperref[eq:6.7]{(7)} is a meromorphic function with periods 2K and 2iK'. Such functions are known as elliptic functions. The connection between elliptic integrals
% Original PDF physical page 257
\sourcepage{257}
and elliptic functions was discovered, but not published, by Gauss; it was rediscovered by Abel and Jacobi\index{Jacobi K G J@Jacobi, K. G. J.}.

\subsection*{EXERCISES}
\originalnumber{1}{exercise:6.2.3.3.1}
1. Prove that formula \hyperref[eq:6.7]{(7)} gives $F(\infty) = iK'$.

\originalnumber{2}{exercise:6.2.3.3.2}
2. Show that $K = K'$ if and only if $k = (\sqrt{2} - 1)^2$.

\originalnumber{3}{exercise:6.2.3.3.3}
3. Show that $f(z), f(z + K)$, and $f(z + iK')$ are odd functions of $z$ while $f(z + K / 2)$ and $f(z + K / 2 + iK')$ are even.

\setcounter{section}{2}\setcounter{subsection}{3}\subsection{The Triangle Functions of Schwarz}\label{sec:6.2.4}
The upper half plane is mapped on a triangle with angles $\alpha_{1}\pi$, $\alpha_{2}\pi$, $\alpha_{3}\pi$ by
\begin{equation*}
F (w) = \int_ {0} ^ {w} w ^ {\alpha_ {1} - 1} (w - 1) ^ {\alpha_ {2} - 1} d w.
\end{equation*}
There are no accessory parameters, as we have already noted.

The inverse function $f(z)$ can again be extended to neighboring triangles by reflection over the sides. This process is particularly interesting when it leads, as in the case of a rectangle, to a meromorphic function. In order that this be so it is necessary that repeated reflections across sides with a common end point should ultimately lead back to the original triangle in an even number of steps. In other words, the angles must be of the form $\pi/n_{1}, \pi/n_{2}, \pi/n_{3}$ with integral denominators. Elementary reasoning shows that the condition
\begin{equation*}
\frac {1}{n _ {1}} + \frac {1}{n _ {2}} + \frac {1}{n _ {3}} = 1
\end{equation*}
is fulfilled only by the triples $(3,3,3)$, $(2,4,4)$, and $(2,3,6)$. They correspond to an equilateral triangle, an isosceles right triangle, and half an equilateral triangle.

In each case it is easy to verify that the reflected images of the triangle fill out the plane, without overlapping and without gaps. This shows that the mapping functions are indeed restrictions of meromorphic functions, known as the Schwarz triangle functions\index{Schwarz triangle function}\index{Triangle function}.

The reader is urged to draw a picture of the triangle net in each of the three cases. He will then observe that each triangle function has a pair of periods with nonreal ratio, and is thus an elliptic function. As an exercise, the reader should determine how many triangles there are in a parallelogram spanned by the periods.

\setcounter{section}{2}\section{A Closer Look At Harmonic Functions}\label{sec:6.3}
We have already discussed the basic properties of harmonic functions\index{Function!harmonic|(}\index{Harmonic function|(}\index{Function!harmonic}\index{Harmonic function} in Chap.~\hyperref[chap:4]{4}, Sec.~\hyperref[sec:4.6]{6}. At that time it was expedient to use a rather crude
% Original PDF physical page 258
\index{Polygon!conformal mapping of|)}\index{Rectangle, mapping on a|)}\sourcepage{258}
definition, namely one that requires all second-order derivatives to be continuous. This was sufficient to prove the mean-value property\index{Mean-value property|(}\index{Mean-value property} from which we could in turn derive the Poisson representation and the reflection principle. We shall now show that a more satisfactory theory is obtained if we make the mean-value property rather than the Laplace equation our starting point.

In this connection we shall also derive an important theorem on monotone sequences of harmonic functions\index{Function!harmonic}\index{Harmonic function}, usually referred to as Harnack's principle.

\setcounter{section}{3}\setcounter{subsection}{0}\subsection{Functions with the Mean-value Property}\label{sec:6.3.1}
Let $u(z)$ be a real-valued continuous function in a region $\Omega$. We say that u satisfies the mean-value property if
\begin{equation*}
u (z _ {0}) = \frac {1}{2 \pi} \int_ {0} ^ {2 \pi} u (z _ {0} + r e ^ {i \theta}) d \theta\tag{8}\label{eq:6.8}
\end{equation*}
when the disk $|z - z_{0}| \leq r$ is contained in $\Omega$. We showed in Chap.~\hyperref[chap:4]{4} that the mean-value property implies the maximum principle. Actually, closer examination of the proof shows that it is sufficient to assume that \hyperref[eq:6.8]{(8)} holds for sufficiently small $r, r < r_{0}$, where we may even allow $r_{0}$ to depend on $z_{0}$. We repeat the conclusion: a continuous function with this property cannot have a relative maximum (or minimum) without reducing to a constant.

We have shown earlier that every harmonic function satisfies the mean-value condition, and we shall now prove the following converse:

\begin{originalstatement}{Theorem}{6}\label{thm:6.6}
A continuous function $u(z)$ which satisfies condition \hyperref[eq:6.8]{(8)} is necessarily harmonic.
\end{originalstatement}
Again, the condition need be satisfied only for sufficiently small $r$. If $u$ satisfies \hyperref[eq:6.8]{(8)}, so does the difference between $u$ and any harmonic function. Suppose that the disk $|z - z_0| \leq \rho$ is contained in $\Omega$, the region where $u$ is defined. By use of Poisson's formula (Chap.~\hyperref[chap:4]{4}, Sec.~\hyperref[sec:4.6.3]{6.3}) we can construct a function $v(z)$ which is harmonic for $|z - z_0| < \rho$, continuous and equal to $u(z)$ on $|z - z_0| = \rho$. The maximum and minimum principle, applied to $u - v$, implies that $u(z) = v(z)$ in the whole disk, and consequently $u(z)$ is harmonic.

The implication of Theorem~\hyperref[thm:6.6]{6} is that we may, if we choose, define a harmonic function to be a continuous function with the mean-value property. Such a function has automatically continuous derivatives of all orders, and it satisfies Laplace's equation.

An analogous reasoning shows that even without the condition \hyperref[eq:6.8]{(8)}


% Original PDF physical page 259
\sourcepage{259}
the assumptions about the derivatives can be relaxed to a considerable degree. Suppose merely that $u(z)$ is continuous and that the derivatives $\partial^{2}u/\partial x^{2}$, $\partial^{2}u/\partial y^{2}$ exist and satisfy $\Delta u = 0$. With the same notations as above we show that the function
\begin{equation*}
V = u - v + \varepsilon (x - x _ {0}) ^ {2},
\end{equation*}
$\varepsilon > 0$, must obey the maximum principle. Indeed, if $V$ had a maximum the rules of the calculus would yield $\partial^{2}V/\partial x^{2} \leq 0$, $\partial^{2}V/\partial y^{2} \leq 0$, and hence $\Delta V \leq 0$ at that point. On the other hand,
\begin{equation*}
\Delta V = \Delta u - \Delta v + 2 \varepsilon = 2 \varepsilon > 0.
\end{equation*}
The contradiction shows that the maximum principle obtains. We can thus conclude that $u - v + \varepsilon(x - x_{0})^{2} \leq \varepsilon\rho^{2}$ in the disk $|z - z_{0}| \leq \rho$. Letting $\varepsilon$ tend to zero we find $u \leq v$, and the opposite inequality can be proved in the same way. Hence u is harmonic.\hyperlink{fn:259:dagger}{\textsuperscript{\ensuremath{\dagger}}}

\setcounter{section}{3}\setcounter{subsection}{1}\subsection{Harnack's Principle}\index{Harnack's principle|(}\index{Harnack's principle}\label{sec:6.3.2}
We recall that Poisson's formula (Chap.~\hyperref[chap:4]{4}, Sec.~\hyperref[sec:4.6.3]{6.3}) permits us to express a harmonic function\index{Function!harmonic}\index{Harmonic function} through its values on a circle. To fit our present needs we write it in the form
\begin{equation*}
u (z) = \frac {1}{2 \pi} \int_ {0} ^ {2 \pi} \frac {\rho^ {2} - r ^ {2}}{| \rho e ^ {i \theta} - z | ^ {2}} u (\rho e ^ {i \theta}) d \theta\tag{9}\label{eq:6.9}
\end{equation*}
where $|z| = r < \rho$ and u is assumed to be harmonic in $|z| \leq \rho$ (or harmonic for $|z| < \rho$, continuous for $|z| \leq \rho$ ). Together with the second of the elementary inequalities
\begin{equation*}
\frac {\rho - r}{\rho + r} \leq \frac {\rho^ {2} - r ^ {2}}{| \rho e ^ {i \theta} - z | ^ {2}} \leq \frac {\rho + r}{\rho - r}\tag{10}\label{eq:6.10}
\end{equation*}
formula \hyperref[eq:6.9]{(9)} yields the estimate
\begin{equation*}
| u (z) | \leq \frac {1}{2 \pi} \frac {\rho + r}{\rho - r} \int_ {0} ^ {2 \pi} | u (\rho e ^ {i \theta}) | d \theta .
\end{equation*}
If it is known that $u(\rho e^{i\theta}) \geq 0$ we can use the first inequality \hyperref[eq:6.10]{(10)} as well, and obtain a double estimate
\begin{equation*}
\frac {1}{2 \pi} \frac {\rho - r}{\rho + r} \int_ {0} ^ {2 \pi} u d \theta \leq u (z) \leq \frac {1}{2 \pi} \frac {\rho + r}{\rho - r} \int_ {0} ^ {2 \pi} u d \theta .
\end{equation*}
But the arithmetic mean of $u(\rho e^{i\theta})$ equals $u(0)$, and we end up with the following upper and lower bounds:
\begin{equation*}
\frac {\rho - r}{\rho + r} u (0) \leq u (z) \leq \frac {\rho + r}{\rho - r} u (0).\tag{11}\label{eq:6.11}
\end{equation*}
\begingroup\renewcommand{\thefootnote}{}\footnotetext{\hypertarget{fn:259:dagger}{}\textsuperscript{\ensuremath{\dagger}} This proof is due to C. Carathéodory\index{Caratheodory C@Caratheodory, C.|hyperpagenote}.}\endgroup

% Original PDF physical page 260
\index{Mean-value property|)}\sourcepage{260}
This is Harnack's inequality\index{Harnack's inequality}; we emphasize that it is valid only for positive harmonic functions\index{Function!harmonic}\index{Harmonic function}.

The main application of \hyperref[eq:6.11]{(11)} is to series with positive terms or, equivalently, increasing sequences of harmonic functions. It leads to a powerful and simple theorem known as Harnack's principle\index{Harnack's principle}:

\begin{originalstatement}{Theorem}{7}\label{thm:6.7}
Consider a sequence of functions $u_{n}(z)$, each defined and harmonic in a certain region $\Omega_{n}$. Let $\Omega$ be a region such that every point in $\Omega$ has a neighborhood contained in all but a finite number of the $\Omega_{n}$, and assume moreover that in this neighborhood $u_{n}(z) \leq u_{n+1}(z)$ as soon as $n$ is sufficiently large. Then there are only two possibilities: either $u_{n}(z)$ tends uniformly to $+\infty$ on every compact subset of $\Omega$, or $u_{n}(z)$ tends to a harmonic limit function $u(z)$ in $\Omega$, uniformly on compact sets.
\end{originalstatement}
The simplest situation occurs when the functions $u_{n}(z)$ are harmonic and form a nondecreasing sequence in $\Omega$. There are, however, applications for which this case is not sufficiently general.

For the proof, suppose first that $\lim_{n\to \infty}u_n(z_0) = \infty$ for at least one point $z_0\in \Omega$. By assumption there exist $r$ and $m$ such that the functions $u_{n}(z)$ are harmonic and form a nondecreasing sequence for $|z - z_0| < r$ and $n\geq m$. If the left-hand inequality \hyperref[eq:6.11]{(11)} is applied to the nonnegative functions $u_{n} - u_{m}$, it follows that $u_{n}(z)$ tends uniformly to $\infty$ in the disk $|z - z_0|\leq r / 2$. On the other hand, if $\lim_{n\to \infty}u_n(z_0) < \infty$, application of the right-hand inequality shows in the same way that $u_{n}(z)$ is bounded on $|z - z_0|\leq r / 2$. Therefore the sets on which $\lim u_n(z)$ is, respectively, finite or infinite are both open, and since $\Omega$ is connected, one of the sets must be empty. As soon as the limit is infinite at a single point, it is hence identically infinite. The uniformity follows by the usual compactness argument.

In the opposite case the limit function $u(z)$ is finite everywhere. With the same notations as above $u_{n+p}(z) - u_n(z) \leq 3(u_{n+p}(z_0) - u_n(z_0))$ for $|z - z_0| \leq r/2$ and $n + p \geq n \geq m$. Hence convergence at $z_0$ implies uniform convergence in a neighborhood of $z_0$, and use of the Heine-Borel property shows that the convergence is uniform on every compact set. The harmonicity of the limit function can be inferred from the fact that $u(z)$ can be represented by Poisson's formula.

\subsection*{EXERCISES}
\originalnumber{1}{exercise:6.3.2.4.1}
1. If $E$ is a compact set in a region $\Omega$, prove that there exists a constant $M$, depending only on $E$ and $\Omega$, such that every positive harmonic function $u(z)$ in $\Omega$ satisfies $u(z_{2}) \leq Mu(z_{1})$ for any two points $z_{1}, z_{2} \in E$.


% Original PDF physical page 261
\index{Function!harmonic|)}\index{Harmonic function|)}\index{Harnack's principle|)}\sourcepage{261}
\setcounter{section}{3}\section{The Dirichlet Problem}\index{Dirichlet's problem|(}\label{sec:6.4}
The most important problem in the theory of harmonic functions is that of finding a harmonic function with given boundary values; it is known as the Dirichlet problem. Poisson's formula solves the problem for a disk, but the case of an arbitrary region is much more difficult. Many methods of solution are known, but none as simple and as suitable for presentation in an elementary text as the method of O. Perron\index{Perron O@Perron, O.}, which is based on the use of subharmonic functions\index{Subharmonic function}.

\setcounter{section}{4}\setcounter{subsection}{0}\subsection{Subharmonic Functions}\label{sec:6.4.1}
Laplace's equation in one dimension would have the form $d^{2}u/dx^{2}=0$. The harmonic functions of one variable would thus be the linear functions $u=ax+b$. A function $v(x)$ is said to be convex if, in any interval, it is at most equal to the linear function $u(x)$ with the same values as $v(x)$ at the end points of the interval.

If this situation is generalized to two dimensions, we are led to the class of subharmonic functions. Linear functions correspond to harmonic functions, intervals correspond to regions, and the end points of an interval correspond to the boundary of the region. Accordingly, a function $v(z)$ of one complex or two real variables will be called subharmonic if in any region $v(z)$ is less than or equal to the harmonic function $u(z)$ which coincides with $v(z)$ on the boundary of the region. Since this formulation requires that we can solve the Dirichlet problem it is preferable to replace the condition by the simpler requirement that $v(z) \leq u(z)$ on the boundary of the region implies $v(z) \leq u(z)$ in the region.

An equivalent but in some respects simpler formulation is the following:

\begin{originalstatement}{Definition}{1}\label{def:6.1}
A continuous real-valued function $v(z)$, defined in a region $\Omega$, is said to be subharmonic in $\Omega$ if for any harmonic function $u(z)$ in a region $\Omega' \subset \Omega$ the difference $v - u$ satisfies the maximum principle in $\Omega'$.
\end{originalstatement}
The condition means that v - u cannot have a maximum in $\Omega'$ without being identically constant. In particular, v itself can have no maximum in $\Omega$. It is important to note that the definition has local character: if v is subharmonic in a neighborhood of each point $z \in \Omega$, then it is subharmonic in $\Omega$. The proof is immediate. A function is said to be subharmonic at a point $z_{0}$ if it is subharmonic in a neighborhood of $z_{0}$. Hence a function is subharmonic in a region if and only if it is subharmonic at all points of the region.

A harmonic function is trivially subharmonic.

A sufficient condition for subharmonicity is that v has a positive Laplacian\index{Laplacian}. In fact, if v - u has a maximum it follows by elementary
% Original PDF physical page 262
\sourcepage{262}
calculus that $\partial^{2}/\partial x^{2}(v-u)\leq0,\partial^{2}/\partial y^{2}(v-u)\leq0$ at that point, provided that these second derivatives exist; this would imply $\Delta v=\Delta(v-u)\leq0$. The condition is not necessary, and as a matter of fact a subharmonic function need not have partial derivatives. If the function has continuous derivatives of the first and second order, it can be shown that the condition $\Delta v\geq0$ is necessary and sufficient. Since we shall not need this property, its proof will be relegated to the exercise section. The condition yields a simple way to ascertain whether a given elementary function of x and y is subharmonic.

We show now that subharmonic functions can be characterized by an inequality which generalizes the mean-value property of harmonic functions:

\begin{originalstatement}{Theorem}{8}\label{thm:6.8}
A continuous function $v(z)$ is subharmonic in $\Omega$ if and only if it satisfies the inequality
\begin{equation*}
v (z _ {0}) \leq \frac {1}{2 \pi} \int_ {0} ^ {2 \pi} v (z _ {0} + r e ^ {i \theta}) d \theta\tag{12}\label{eq:6.12}
\end{equation*}
for every disk $|z - z_0| \leq r$ contained in $\Omega$.

\end{originalstatement}
The sufficiency follows by the fact that \hyperref[eq:6.12]{(12)}, rather than the mean-value property, is what is actually needed in order to show that v cannot have a maximum without being constant. Since v - u satisfies the same inequality, it follows that v is subharmonic.

In order to prove the necessity we form the Poisson integral $P_{v}(z)$ in the disk $|z - z_{0}| < r$ with the values of v taken on the circumference $|z - z_{0}| = r$. If v is subharmonic, the function $v - P_{v}$ can have no maximum in the disk unless it is constant. By Schwarz's theorem (Chap.~\hyperref[chap:4]{4}, Theorem 25) $v - P_{v}$ tends to 0 as z approaches a point on the circumference. Hence $v - P_{v}$ has a maximum in the closed disk. If the maximum were positive it would be taken at an interior point, and the function could not be constant. This is a contradiction, and we conclude that $v \leq P_{v}$. For $z = z_{0}$ we obtain $v(z_{0}) \leq P_{v}(z_{0})$, and this is the inequality \hyperref[eq:6.12]{(12)}.

We list now a number of elementary properties of subharmonic functions:

\textit{1. If $v$ is subharmonic, so is $kv$ for any constant $k \geq 0$.}

\textit{2. If $v_{1}$ and $v_{2}$ are subharmonic, so is $v_{1} + v_{2}$.}

These are immediate consequences of Theorem~\hyperref[thm:6.8]{8}. The next property follows most easily from the original definition.

\textit{3. If $v_{1}$ and $v_{2}$ are subharmonic in $\Omega$, then $v = \max (v_{1}, v_{2})$ is likewise subharmonic in $\Omega$.}


% Original PDF physical page 263
\sourcepage{263}
The notation is to be understood in the sense that $v(z)$ is at each point equal to the greater of the values $v_{1}(z)$ and $v_{2}(z)$. The continuity of v is obvious. Suppose now that v - u has a maximum at $z_{0} \in \Omega'$ where u is defined and harmonic in $\Omega'$. We may assume that $v(z_{0}) = v_{1}(z_{0})$. Then
\begin{equation*}
v _ {1} (z) - u (z) \leq v (z) - u (z) \leq v (z _ {0}) - u (z _ {0}) = v _ {1} (z _ {0}) - u (z _ {0})
\end{equation*}
for $z \in \Omega'$. Hence $v_1 - u$ is constant, and by the same inequality $v - u$ must also be constant. It is proved that $v$ is subharmonic.

Let $\Delta$ be a disk whose closure is contained in $\Omega$, and denote by $P_{v}$ the Poisson integral formed with the values of v on its circumference. Then the following is true:

\textit{4. If $v$ is subharmonic, then the function $v'$ defined as $P_v$ in $\Delta$ and as $v$ outside of $\Delta$ is also subharmonic.}

The continuity of $v'$ follows by the theorem of Schwarz. We have proved that $v \leq P_{v}$ in $\Delta$, and hence $v \leq v'$ throughout $\Omega$. It is clear that $v'$ is subharmonic in the interior and exterior of $\Delta$. Suppose now that $v' - u$ had a maximum at a point $z_{0}$ on the circumference of $\Delta$. It follows at once that v - u would also have a maximum at $z_{0}$. Hence v - u would be constant, and the inequality
\begin{equation*}
v - u \leq v ^ {\prime} - u \leq v ^ {\prime} (z _ {0}) - u (z _ {0}) = v (z _ {0}) - u (z _ {0})
\end{equation*}
shows that $v' - u$ is likewise constant. We conclude that $v'$ is subharmonic.

\emph{Remark.} We are considering only continuous subharmonic functions, but the generally accepted definition requires merely that the function be upper semicontinuous. A real-valued function $v(z)$ is upper semicontinuous (u.s.c.) at $z_0$ if $\limsup_{z \to z_0} v(z) \leq v(z_0)$ and lower semicontinuous (l.s.c.) if $\liminf_{z \to z_0}$

$v(z) \geq v(z_{0})$. If in doubt, which is which, remember that upper refers to the upper half and lower to the lower half of the double inequality $v(z_{0}) - \varepsilon < v(z) < v(z_{0}) + \varepsilon$. It is also customary to allow an u.s.c. function to assume the value $-\infty$ and a l.s.c. function the value $+\infty$.

In all other respects Definition~\hyperref[def:6.1]{1} is unchanged. The maximum principle is as meaningful for upper semicontinuous as for continuous functions due to the fact that an upper semicontinuous function\index{Function!semicontinuous} will also attain a maximum on any compact set (see Ex.~\hyperref[exercise:6.4.1.5.6]{6}).

It can also be shown that the integral in \hyperref[eq:6.12]{(12)} has a meaning and that Theorem~\hyperref[thm:6.8]{8} remains valid when v is only u.s.c.

\subsection*{EXERCISES}
\originalnumber{1}{exercise:6.4.1.5.1}
1. Show that the functions $|x|$, $|z|^{\alpha}(\alpha \geq 0)$, $\log(1 + |z|^{2})$ are subharmonic.


% Original PDF physical page 264
\sourcepage{264}
\originalnumber{2}{exercise:6.4.1.5.2}
2. If $f(z)$ is analytic, prove that $|f(z)|^{\alpha} (\alpha \geq 0)$ and $\log (1 + |f(z)|^{2})$ are subharmonic.

\originalnumber{3}{exercise:6.4.1.5.3}
3. If $v$ is continuous together with its partial derivatives up to the second order, prove that $v$ is subharmonic if and only if $\Delta v \geq 0$. Hint: For the sufficiency, prove first that $v + \varepsilon x^{2}$, $\varepsilon > 0$, is subharmonic. For the necessity, show that if $\Delta v < 0$ the mean value over a circle would be a decreasing function of the radius.

\originalnumber{4}{exercise:6.4.1.5.4}
4. Prove that a subharmonic function remains subharmonic if the independent variable is subjected to a conformal mapping\index{Conformal mapping}\index{Mapping!conformal}.

\originalnumber{5}{exercise:6.4.1.5.5}
5. Formulate and prove a theorem to the effect that a uniform limit of subharmonic functions is subharmonic.

\originalnumber{6}{exercise:6.4.1.5.6}
6. If $v(z)$ is upper semicontinuous on the open set $\Omega$, show that it has a maximum on any compact set $E \subset \Omega$.

\setcounter{section}{4}\setcounter{subsection}{1}\subsection{Solution of Dirichlet's Problem}\label{sec:6.4.2}
The first to use subharmonic functions for the study of Dirichlet's problem\index{Dirichlet's problem} was O. Perron\index{Perron O@Perron, O.}. His method is characterized by extreme generality, and it is completely elementary.

We consider a bounded region $\Omega$ and a real-valued function $f(\zeta)$ defined on its boundary $\Gamma$ (for clarity, boundary points will be denoted by $\zeta$). To begin with, $f(\zeta)$ need not even be continuous, but for the sake of simplicity we assume that it is bounded, $|f(\zeta)| \leq M$. With each $f$ we associate a harmonic function $u(z)$ in $\Omega$, defined by a simple process which will be detailed below. If $f$ is continuous, and if $\Omega$ satisfies certain mild conditions, the corresponding function $u$ will solve the Dirichlet problem for $\Omega$ with the boundary values $f$.

We define the class $\mathfrak{B}(f)$ of functions $v$ with the following properties:

(a) $v$ is subharmonic in $\Omega$;

(b) $\varlimsup_{z\to \zeta}v(z)\leq f(\zeta)$ for all $\zeta \in \Gamma$

The precise meaning of (b) is this: given $\varepsilon > 0$ and a point $\zeta \in \Gamma$ there exists a neighborhood $\Delta$ of $\zeta$ such that $v(z) < f(\zeta) + \varepsilon$ in $\Delta \cap \Omega$. The class $\mathfrak{B}(f)$ is not empty, for it contains all constants $\leq -M$. We prove:

\begin{originalstatement}{Lemma}{1}\label{lem:6.1}
The function $u$, defined as $u(z) = \text{l.u.b. } v(z)$ for $v \in \mathfrak{B}(f)$, is harmonic in $\Omega$.
\end{originalstatement}
In the first place, each $v$ is $\leq M$ in $\Omega$. This is a simple enough consequence of the maximum principle, but because of its importance we want to explain this point in some detail. For a given $\varepsilon > 0$, let $E$ be the set of points $z \in \Omega$ for which $v(z) \geq M + \varepsilon$. The points $z$ in the complement
% Original PDF physical page 265
\sourcepage{265}
$\sim E$ are of three kinds: (1) points in the exterior of $\Omega$, (2) points on $\Gamma$, (3) points in $\Omega$ with $v(z) < M + \varepsilon$. In case (1) $z$ has a neighborhood contained in the exterior, in case (2) there is a neighborhood $\Delta$ with $v < M + \varepsilon$ in $\Delta \cap \Omega$, by property (b), and in case (3) there exists, by continuity, a neighborhood in $\Omega$ with $v < M + \varepsilon$. Hence $\sim E$ is open, and $E$ is closed. Moreover, since $\Omega$ is bounded, $E$ is compact. If $E$ were not void, $v$ would have a maximum on $E$, and this would also be a maximum in $\Omega$. This is impossible, for because of (b) $v$ cannot be a constant $> M$. Hence $E$ is void for every $\varepsilon$, and it follows that $v \leq M$ in $\Omega$.

Consider a disk $\Delta$ whose closure is contained in $\Omega$, and a point $z_0 \in \Delta$. There exists a sequence of functions $v_n \in \mathfrak{B}(f)$ such that $\lim_{n \to \infty} v_n(z_0) = u(z_0)$.

Set $V_{n} = \max (v_{1}, v_{2}, \ldots, v_{n})$. Then the $V_{n}$ form a nondecreasing sequence of functions in $\mathfrak{B}(f)$. We construct $V_{n}'$ equal to $V_{n}$ outside of $\Delta$ and equal to the Poisson integral of $V_{n}$ in $\Delta$. By property \hyperref[eq:6.4]{(4)} of the preceding section the $V_{n}'$ are still in $\mathfrak{B}(f)$. They form a nondecreasing sequence, and the inequality $v_{n}(z_{0}) \leq V_{n}(z_{0}) \leq V_{n}'(z_{0}) \leq u(z_{0})$ shows that $\lim_{n \to \infty} V_{n}'(z_{0}) = u(z_{0})$. By Harnack's principle the sequence $\{V_{n}'\}$ converges to a harmonic limit function $U$ in $\Delta$ which satisfies $U \leq u$ and $U(z_{0}) = u(z_{0})$.

Suppose now that we start the same process from another point $z_1 \in \Delta$. We select $w_n \in \mathfrak{B}(f)$ so that $\lim_{n \to \infty} w_n(z_1) = u(z_1)$, but this time, before proceeding with the construction, we replace $w_n$ by $\bar{w}_n = \max (v_n, w_n)$. Setting $W_n = \max (\bar{w}_1, \ldots, \bar{w}_n)$ we construct the corresponding sequence $\{W_n'\}$ with the aid of the Poisson integral and are led to a harmonic limit function $U_1$ which satisfies $U \leq U_1 \leq u$ and $U_1(z_1) = u(z_1)$. It follows that $U - U_1$ has the maximum zero at $z_0$. Therefore $U$ is identically equal to $U_1$, and we have proved that $u(z_1) = U(z_1)$ for arbitrary $z_1 \in \Delta$. It follows that $u$ is harmonic in any disk $\Delta$ and, consequently, in all of $\Omega$.

We will now investigate the circumstances under which u solves the Dirichlet problem for continuous f. We note first that the Dirichlet problem does not always have a solution. For instance, if $\Omega$ is the punctured disk $0 < |z| < 1$, consider the boundary values $f(0) = 1$ and $f(\zeta) = 0$ for $|\zeta| = 1$. A harmonic function with these boundary values would be bounded and would, hence, present a removable singularity at the origin. But then the maximum principle would imply that the function vanishes identically and thus could not have the boundary value 1 at the origin. It follows that no solution can exist.

It is also easy to see that a solution, if it exists, must be identical with u. In fact, if U is a solution it is first of all clear that $U \in \mathfrak{B}(f)$, and hence $u \geq U$. The opposite inequality $u \leq U$ follows by the maximum principle which implies $v \leq U$ for all $v \in \mathfrak{B}(f)$.


% Original PDF physical page 266
\sourcepage{266}
The existence of a solution can be asserted for a wide class of regions. Generally speaking, the solution exists if the complement of $\Omega$ is not too “thin” in the neighborhood of any boundary point. We begin by proving a lemma which, on the surface, seems to have little to do with the notion of thinness.

\begin{originalstatement}{Lemma}{2}\label{lem:6.2}
Suppose that there exists a harmonic function $\omega(z)$ in $\Omega$ whose continuous boundary values $\omega(\zeta)$ are strictly positive except at one point $\zeta_0$ where $\omega(\zeta_0) = 0$. Then, if $f(\zeta)$ is continuous at $\zeta_0$, the corresponding function $u$ determined by Perron's method satisfies $\lim u(z) = f(\zeta_0)$.
\end{originalstatement}
The lemma will be proved if we show that $\varlimsup_{z\to \zeta_0}u(z)\leq f(\zeta_0) + \varepsilon$ and $\varliminf_{z\to \zeta_0}u(z)\geq f(\zeta_0) - \varepsilon$ for all $\varepsilon >0$. We are still assuming that $\Omega$ is bounded and $|f(\zeta)|\leq M$.

Determine a neighborhood $\Delta$ of $\zeta_0$ such that $|f(\zeta) - f(\zeta_0)| < \varepsilon$ for $\zeta \in \Delta$. In $\Omega - \Delta \cap \Omega$ the function $\omega(z)$ has a positive minimum $\omega_0$. We consider the boundary values of the harmonic function
\begin{equation*}
W (z) = f (\zeta_ {0}) + \varepsilon + \frac {\omega (z)}{\omega_ {0}} (M - f (\zeta_ {0})).
\end{equation*}
For $\zeta \in \Delta$ we have $W(\zeta) \geq f(\zeta_0) + \varepsilon > f(\zeta)$, and for $\zeta$ outside of $\Delta$ we obtain $W(\zeta) \geq M + \varepsilon > f(\zeta)$. By the maximum principle any function $v \in \mathfrak{B}(f)$ must hence satisfy $v(z) < W(z)$. It follows that $u(z) \leq W(z)$ and consequently $\varlimsup_{z \to \zeta_0} u(z) \leq W(\zeta_0) = f(\zeta_0) + \varepsilon$, which is the first inequality we set out to prove.

For the second inequality we need only show that the function
\begin{equation*}
V (z) = f (\zeta_ {0}) - \varepsilon - \frac {\omega (z)}{\omega_ {0}} (M + f (\zeta_ {0}))
\end{equation*}
is in $\mathfrak{B}(f)$. For $\zeta \in \Delta$ we have $V(\zeta) \leq f(\zeta_0) - \varepsilon < f(\zeta)$, and at all other boundary points $V(\zeta) \leq -M - \varepsilon < f(\zeta)$. Since $V$ is harmonic it belongs to $\mathfrak{B}(f)$ and we obtain $u(z) \geq V(z)$, $\underline{\lim}_{z \to \zeta_0} u(z) \geq V(\zeta_0) = f(\zeta_0) - \varepsilon$. This completes the proof.

The function $\omega(z)$ of Lemma~\hyperref[lem:6.2]{2} is sometimes called a barrier\index{Barrier} at the point $\zeta_{0}$. Clearly, we can now say that the Dirichlet problem is solvable provided that there is a barrier at each boundary point. It remains to formulate geometric conditions which imply the existence of a barrier. Necessary and sufficient conditions are known, but they are not purely geometric, and therefore difficult to apply. It is relatively easy, however, to find sufficient conditions with a wide range of applicability.


% Original PDF physical page 267
\sourcepage{267}
To\index{Conformal equivalence|(} begin with the simplest case, suppose that $\Omega \cup \Gamma$ is contained in an open half plane, except for a point $\zeta_0$ which lies on the boundary line. If the direction of this line is $\alpha$ (with the half plane to the left), then $\omega(z) = \operatorname{Im} e^{-i\alpha}(z - \zeta_0)$ is a barrier at $\zeta_0$.

More generally, suppose that $\zeta_{0}$ is the end point of a line segment all of whose points, except $\zeta_{0}$, lie in the exterior of $\Omega$. If the other end point is denoted by $\zeta_{1}$, we know that a single-valued branch of
\begin{equation*}
\sqrt {\frac {z - \zeta_ {0}}{z - \zeta_ {1}}}
\end{equation*}
can be defined outside of the segment. With a proper determination of the angle $\alpha$ the function
\begin{equation*}
\operatorname{Im} \left[ e ^ {- i \alpha} \sqrt {\frac {z - \zeta_ {0}}{z - \zeta_ {1}}} \right]
\end{equation*}
is easily seen to be a barrier at $\zeta_{0}$.

This is not the strongest result that can be obtained by these methods, but it is sufficient for most applications. We shall therefore be content with the following statement:

\begin{originalstatement}{Theorem}{9}\label{thm:6.9}
The Dirichlet problem can be solved for any region $\Omega$ such that each boundary point is the end point of a line segment whose other points are exterior to $\Omega$.
\end{originalstatement}
The hypothesis is fulfilled if $\Omega$ and its complement have a common boundary consisting of a finite number of simple closed curves with a tangent at each point. Corners and certain types of cusps are also permissible. \hyperlink{fn:267:dagger}{\textsuperscript{\ensuremath{\dagger}}}

\subsection*{EXERCISE}
\originalnumber{}{exercise:6.4.2.6.unnumbered}
If $\Omega$ is the punctured disk $0 < |z| < 1$ and if $f$ is given by $f(\zeta) = 0$ for $|\zeta| = 1$, $f(0) = 1$, show that all functions $v \in \mathfrak{B}(f)$ are $\leq 0$ in $\Omega$.

\setcounter{section}{4}\section{Canonical Mappings Of Multiply Connected Regions}\index{Canonical mapping|(}\index{Canonical mapping}\label{sec:6.5}
Riemann's mapping theorem permits us to conclude that any two simply connected regions, with the exception of the whole plane, can be mapped conformally onto each other, or that they are conformally equivalent.

\begingroup\renewcommand{\thefootnote}{}\footnotetext{\hypertarget{fn:267:dagger}{}\textsuperscript{\ensuremath{\dagger}} The best result that can be proved by essentially the same method is the following: \emph{The Dirichlet problem can be solved for any region whose complement is such that no component reduces to a point.} From this proposition an independent proof of the Riemann mapping theorem can easily be derived.}\endgroup

% Original PDF physical page 268
\index{Conformal equivalence|)}\index{Dirichlet's problem|)}\sourcepage{268}
For multiply connected regions of the same connectivity this is no longer true. Instead we must try to find a system of canonical regions\index{Canonical region} with the property that each multiply connected region is conformally equivalent to one and only one canonical region. The choice of canonical regions is to a certain extent arbitrary, and there are several types with equally simple properties.

In order to stay on an elementary level we will limit ourselves to the study of regions of finite connectivity. We shall find that the basic step toward the construction of canonical mappings\index{Canonical mapping} is the introduction of certain harmonic functions with a particularly simple behavior on the boundary. Of these the harmonic measures are related only to the region and one of its contours, while the Green's function\index{Function!Green's} is related to the region and an interior point.

\setcounter{section}{5}\setcounter{subsection}{0}\subsection{Harmonic Measures}\label{sec:6.5.1}
When studying the conformal mappings\index{Conformal mapping}\index{Mapping!conformal} of a region $\Omega$ we can of course replace $\Omega$ by any region known to be conformally equivalent to $\Omega$, that is to say, we can perform preliminary conformal mappings at will. Because of this freedom in the choice of the original region it turns out that it is never necessary to deal with the difficulties which may arise from a complicated structure of the boundary.

In the following $\Omega$ denotes a plane region of connectivity n > 1. The components of the complement are denoted by $E_{1}, E_{2}, \ldots, E_{n}$, and we take $E_{n}$ to be the unbounded component. Without loss of generality we can and will assume that no $E_{k}$ reduces to a point, for it is clear that a point component is a removable singularity of any mapping function, and consequently the mappings remain the same if this isolated boundary point is added to the region.

The complement of $E_{n}$ is a simply connected region $\Omega'$. By Riemann's theorem, $\Omega'$ can be mapped conformally onto the disk $|z| < 1$; under this mapping $\Omega$ is transformed into a new region, and the images of $E_{1}, \ldots, E_{n-1}$ are the bounded components of its complement. For the sake of convenience we agree to use the same notations as before the mapping; in particular, $E_{n}$ is now the set $|z| \geq 1$. The unit circle $|z| = 1$, traced in the positive direction, will be denoted by $C_{n}$ and is called the outer contour\index{Contour!outer} of the new region $\Omega$.

Consider now the complement of $E_{1}$ with respect to the extended plane. This is again a simply connected region, and we map it onto the outside of the unit circle with $\infty$ corresponding to itself. The image of $C_{n}$ is a directed closed analytic curve which we continue to denote by $C_{n}$, just as we keep all the other notations. In addition we define the inner contour\index{Contour!inner} $C_{1}$ to be the unit circle in the new plane, traced in the negative direction.


% Original PDF physical page 269
\sourcepage{269}
The process can evidently be repeated until we end up with a region $\Omega$ bounded by an outer contour $C_n$ and $n - 1$ inner contours $C_1, \ldots, C_{n-1}$ (Fig.~\hyperref[fig:6-3]{6-3}). It is important to note that the index of a contour with respect to an arbitrary point in the plane can be readily computed. For instance, at the stage where $C_k$, $k < n$, is the unit circle, the index of $C_k$ is $-1$ with respect to interior points of $E_k$ and $0$ with respect to all other points not on $C_k$. The subsequent mappings will not change this state of affairs. The fact is clear, and a formal proof based on the argument principle can easily be given. One shows in the same way that the outer contour $C_n$ has the index $0$ with respect to interior points of $E_n$ and the index $1$ with respect to all other points not on $C_n$. It follows that the cycle $C = C_1 + C_2 + \cdots + C_n$ bounds $\Omega$ in the sense of Chap.~\hyperref[chap:4]{4}, Sec.~\hyperref[sec:4.5.1]{5.1}, Definition~\hyperref[def:4.4]{4}. The distinction between outer and inner contours is coincidental, for evidently an inversion with respect to an interior point of $E_k$ will make $C_k$ the outer contour.

It is clear that Theorem~\hyperref[thm:6.9]{9} applies to $\Omega$. As a matter of fact the existence of a barrier is completely obvious since any contour can be transformed into a circle.

\bookfigure{269}{1}{FIG. 6-3. Transformations of a multiply connected region.}

% Original PDF physical page 270
\sourcepage{270}
Suppose now that we solve the Dirichlet problem in $\Omega$ with the boundary values 1 on $C_{k}$ and 0 on the other contours. The solution is denoted by $\omega_{k}(z)$, and it is called the harmonic measure of $C_{k}$ with respect to the region $\Omega$. We have clearly $0 < \omega_{k}(z) < 1$ in $\Omega$ and
\begin{equation*}
\omega_ {1} (z) + \omega_ {2} (z) + \dots + \omega_ {n} (z) \equiv 1.
\end{equation*}
If we map $\Omega$ so that $C_{i}$ becomes a circle, then $\omega_{k}$ can be continued across $C_{i}$ according to the reflection principle. We conclude that $\omega_{k}$ is harmonic in the closed region $\Omega$ in the sense that it can be extended to a larger region.

The contours $C_{1}, \ldots, C_{n-1}$ form a homology basis for the cycles in $\Omega$, homology being understood with respect to an unspecified larger region. The conjugate harmonic differential of $\omega_{k}$ has periods
\begin{equation*}
\alpha_ {k j} = \int_ {C _ {j}} \frac {\partial \omega_ {k}}{\partial n} d s = \int_ {C _ {j}} {}^{*}d \omega_ {k}
\end{equation*}
along $C_j$. We assert that no linear combination $\lambda_1\omega_1(z) + \lambda_2\omega_2(z) + \cdots + \lambda_{n-1}\omega_{n-1}(z)$ with constant coefficients can have a single-valued conjugate function unless all the $\lambda_i$ are zero. To see this, suppose that this expression were the real part of an analytic function $f(z)$. By the reflection principle, $f(z)$ would have an analytic extension to the closure of $\Omega$. The real part of $f(z)$ would be constantly equal to $\lambda_i$ on $C_i$, $i = 1, \ldots, n - 1$, and zero on $C_n$. Consequently, each contour would be mapped onto a vertical line segment. If $w_0$ does not lie on any of these segments, a single-valued branch of $\arg (f(z) - w_0)$) can be defined on each contour. It follows by the argument principle that $f(z)$ cannot take the value $w_0$ in $\Omega$. But then $f(z)$ must reduce to a constant, for otherwise the image of $\Omega$ would certainly contain points not on the line segments. We conclude that the real part of $f(z)$ is identically zero, and hence the boundary values $\lambda_i$ must all vanish.

What we have proved is that the homogeneous system of linear equations
\begin{equation*}
\lambda_ {1} \alpha_ {1 j} + \lambda_ {2} \alpha_ {2 j} + \dots + \lambda_ {n - 1} \alpha_ {n - 1, j} = 0 (j = 1, \dots , n - 1)\tag{13}\label{eq:6.13}
\end{equation*}
has only the trivial solution $\lambda_{i}=0$, for these are the conditions under which $\lambda_{1}\omega_{1}+\cdots+\lambda_{n-1}\omega_{n-1}$ has a single-valued conjugate. By the theory of linear equations any inhomogeneous system of equations with the same coefficients as \hyperref[eq:6.13]{(13)} must have a solution. In particular, we conclude that it is possible to solve the system


% Original PDF physical page 271
\sourcepage{271}
\begin{equation*}
\begin{array}{l} \lambda_ {1} \alpha_ {1 1} + \lambda_ {2} \alpha_ {2 1} + \dots + \lambda_ {n - 1} \alpha_ {n - 1, 1} = 2 \pi \\ \lambda_ {1} \alpha_ {1 2} + \lambda_ {2} \alpha_ {2 2} + \dots + \lambda_ {n - 1} \alpha_ {n - 1, 2} = 0 \\ \dots \dots \\ \lambda_ {1} \alpha_ {1, n - 1} + \lambda_ {2} \alpha_ {2, n - 1} + \dots + \lambda_ {n - 1} \alpha_ {n - 1, n - 1} = 0 \\ \lambda_ {1} \alpha_ {1 n} + \lambda_ {2} \alpha_ {2 n} + \dots + \lambda_ {n - 1} \alpha_ {n - 1, n} = - 2 \pi \end{array}\tag{14}\label{eq:6.14}
\end{equation*}
where the last equation is a consequence of the n-1 first (because $\alpha_{k1} + \alpha_{k2} + \cdots + \alpha_{kn} = 0$ ). In other words, we can find a multiple-valued integral $f(z)$ with periods $\pm 2\pi i$ along $C_{1}$ and $C_{n}$ and all other periods equal to zero, the real part being constantly equal to $\lambda_{k}$ on $C_{k}$ (we set $\lambda_{n} = 0$ ). The function $F(z) = e^{f(z)}$ is then single-valued. We prove:

\begin{originalstatement}{Theorem}{10}\label{thm:6.10}
The function $F(z)$ effects a one-to-one conformal mapping\index{Conformal mapping}\index{Mapping!conformal} of $\Omega$ onto the annulus $1 < |w| < e^{\lambda_1}$ minus $n - 2$ concentric arcs situated on the circles $|w| = e^{\lambda_i}$, $i = 2, \ldots, n - 1$.
\end{originalstatement}
The mapping is illustrated in Fig.~\hyperref[fig:6-4]{6-4}. The contours $C_{1}$ and $C_{n}$ are in one-to-one correspondence with the full circles, while the other contours are flattened into circular slits. It should be imagined that each slit has two edges which together with the end points form a closed contour.

The proof is by use of the argument principle. We know that $F(z)$ is analytic with a constant modulus on each contour. The number of roots of the equation $F(z) = w_{0}$ is given by
\begin{equation*}
\begin{array}{r l} \frac {1}{2 \pi i} \int_ {C _ {1}} \frac {F ^ {\prime} (z) d z}{F (z) - w _ {0}} + \frac {1}{2 \pi i} \int_ {C _ {2}} \frac {F ^ {\prime} (z) d z}{F (z) - w _ {0}} + \dots & \\ & + \frac {1}{2 \pi i} \int_ {C _ {n}} \frac {F ^ {\prime} (z) d z}{F (z) - w _ {0}}, \end{array}\tag{15}\label{eq:6.15}
\end{equation*}
at any rate if $w_{0}$ is not taken on the boundary. For $w_{0}=0$ the terms in \hyperref[eq:6.15]{(15)} are known, being equal to 1, 0, . . . , 0, -1, respectively. The
\bookfigure{271}{1}{FIG. 6-4. Concentric slit region.}
% Original PDF physical page 272
\sourcepage{272}
integral\index{Function!Green's|(} over $C_1$ remains constantly equal to 1 for $|w_0| < e^{\lambda_1}$, and it vanishes for $|w_0| > e^{\lambda_1}$; similarly, the last integral is -1 for $|w_0| < 1$ and 0 for $|w_0| > 1$. The integrals over $C_k$, $1 < k < n$, vanish for all $w_0$ with $|w_0| \neq e^{\lambda_k}$. Suppose now that the value $w_0$ is actually taken by $F(z)$; inasmuch as $\Omega$ must be mapped onto an open set, we can choose $|w_0| \neq \text{all } e^{\lambda_i}$. For this $w_0$ the expression \hyperref[eq:6.15]{(15)} must be positive. But that is possible only if $1 < |w_0| < e^{\lambda_1}$. Thus $\lambda_1 > 0$ and, by continuity, $0 \leq \lambda_i \leq \lambda_1$.

From here on the proof could be completed by means of a purely topological argument. It is more instructive, however, and in fact simpler, to draw the conclusion from the argument principle. When there are simple poles on the boundary, the residue theorem continues to hold provided that the contour integral is replaced by its Cauchy principal value, and provided that the sum of the residues includes one-half of the residues on the boundary.\hyperlink{fn:272:dagger}{\textsuperscript{\ensuremath{\dagger}}} In the present situation the second convention means that a value taken on the boundary is counted with half its multiplicity. The computation of the principal values causes no difficulty. If $|w_0| = e^{\lambda_k}$, we find that
\begin{equation*}
\text { pr.v. } \int_ {C _ {k}} \frac {F ^ {\prime} (z) d z}{F (z) - w _ {0}} = \frac {1}{2} \int_ {C _ {k}} \frac {F ^ {\prime} (z) d z}{F (z)},
\end{equation*}
for by elementary geometry (or direct computation)
\begin{equation*}
d \arg (F (z) - w _ {0}) = \frac {1}{2} d \arg F (z).
\end{equation*}
Consequently, the principal values in \hyperref[eq:6.15]{(15)} are $\frac{1}{2}$ for $k = 1$, 0 for $2 \leq k \leq n - 1$, $-\frac{1}{2}$ for $k = n$.

We conclude now that each value on the circle $|w_{0}| = 1$ or $|w_{0}| = e^{\lambda_{1}}$ is taken one-half time, that is to say once on the boundary; this proves that $C_{1}$ and $C_{n}$ are mapped in a one-to-one manner and that $0 < \lambda_{i} < \lambda_{1}, i \neq 1, n$. Next, if $1 < |w_{0}| < e^{\lambda_{1}}$, it follows that $w_{0}$ is taken either once in the interior, twice on the boundary, or once on the boundary with the multiplicity 2. On each contour $C_{2}, \ldots, C_{n-1}$ a single-valued branch of arg $F(z)$ can be defined, and the values of multiplicity 2 correspond to relative maxima and minima of arg $F(z)$. There is at least one maximum and one minimum, and there cannot be more or else $F(z)$ would pass more than twice through the same values. Furthermore, the difference between the maximum and the minimum must be $<2\pi$, which shows that each contour is mapped onto a proper arc. Finally, the arcs which correspond to different contours must be disjoint.

\begingroup\renewcommand{\thefootnote}{}\footnotetext{\hypertarget{fn:272:dagger}{}\textsuperscript{\ensuremath{\dagger}} In Chap.~\hyperref[chap:4]{4}, Sec.~\hyperref[sec:4.5.3]{5.3}, the Cauchy principal value was introduced in the case of an integral over a straight line. In the case of an arbitrary analytic arc it is simplest to define the principal value by means of an auxiliary conformal mapping\index{Conformal mapping}\index{Mapping!conformal} which transforms a subarc into a line segment. The stated generalization of the residue theorem follows quite easily and proves that the principal value is independent of the auxiliary conformal mapping.}\endgroup

% Original PDF physical page 273
\sourcepage{273}
We have proved the complete Theorem~\hyperref[thm:6.10]{10}, and in addition we have been able to describe the correspondence of the boundaries. The significance of the theorem is that we can map $\Omega$ onto a canonical region bounded by two circles and $n - 2$ concentric circular slits; by way of normalization the radius of the inner circle is chosen equal to 1. For a given choice of $C_1$ and $C_n$ the canonical mapping\index{Canonical mapping} is uniquely determined up to a rotation. This follows from the fact that the system \hyperref[eq:6.15]{(15)} has only one solution.

The shape of a canonical region of connectivity n depends on 3n - 6 real constants. In fact, the position and size of each slit is determined by three numbers, a total of 3n - 6; the thickness of the annulus gives one additional parameter, but another parameter must be discounted to allow for the arbitrary rotation.

\subsection*{EXERCISES}
\originalnumber{1}{exercise:6.5.1.7.1}
1. Prove directly that two circular annuli are conformally equivalent if and only if the ratios of their radii are equal.

\originalnumber{2}{exercise:6.5.1.7.2}
2. Prove that $\alpha_{ij} = \alpha_{ji}$. Hint: Apply Theorem~\hyperref[thm:4.21]{21}, Chap.~\hyperref[chap:4]{4}.

\setcounter{section}{5}\setcounter{subsection}{1}\subsection{Green's Function}\label{sec:6.5.2}
We suppose again that $\Omega$ is a region of finite connectivity, and inasmuch as preliminary conformal mappings\index{Conformal mapping}\index{Mapping!conformal} will be permissible we can assume that $\Omega$ is bounded by analytic contours $C_1, \ldots, C_n$; this time the case $n = 1$ will be included.

We consider a point $z_{0} \in \Omega$ and solve the Dirichlet problem in $\Omega$ with the boundary values $\log |\zeta - z_{0}|$. The solution is denoted by $G(z)$, but the main interest is attached to the function $g(z) = G(z) - \log |z - z_{0}|$, known as the Green's function\index{Green's function|(}\index{Function!Green's}\index{Green's function} of $\Omega$ with pole at $z_{0}$. When the dependence on $z_{0}$ is emphasized, it is denoted by $g(z, z_{0})$.

The Green's function is harmonic in $\Omega$ except at $z_0$, and it vanishes on the boundary. In a neighborhood of $z_0$ it differs from $-\log |z - z_0|$ by a harmonic function. By these properties $g(z)$ is uniquely determined. In fact, if $g_1(z)$ has the same properties, then $g - g_1$ is harmonic throughout $\Omega$ and vanishes on the boundary. By the maximum principle it follows that $g_1$ is identically equal to $g$.

If two regions are conformally equivalent, then the Green's functions with corresponding poles are equal at points which correspond to each other. To be more explicit, let $z = z(\zeta)$ define a one-to-one conformal mapping of a region $\Omega'$ in the $\zeta$-plane onto a region $\Omega$ in the $z$-plane. Choose a point $\zeta_0 \in \Omega'$ and denote by $g(z, z_0)$ the Green's function of $\Omega$ with pole at $z_0 = z(\zeta_0)$. It is claimed that $g(z(\zeta), z_0)$ is the Green's function of $\Omega'$. To begin with, if $\zeta$ tends to a boundary point, then $z(\zeta)$ approaches the boundary of $\Omega$, and hence $g(z(\zeta), z_0)$ has the boundary
% Original PDF physical page 274
\sourcepage{274}
values zero. As to the behavior at $\zeta_0$ we know that $g(z(\zeta), z_0)$ differs from $-\log |z(\zeta) - z(\zeta_0)|$ by a harmonic function of $z(\zeta)$, and hence by a harmonic function of $\zeta$. But the difference $\log |z(\zeta) - z(\zeta_0)| - \log |\zeta - \zeta_0|$ is also harmonic, and it follows that $g(z(\zeta), z_0)$ has the desired behavior at $\zeta_0$. We have proved that the Green's function\index{Function!Green's}\index{Green's function} is invariant under conformal mappings\index{Conformal mapping}\index{Mapping!conformal}, and it is in view of this invariance that preliminary conformal mappings can be performed at will.

In the case of a simply connected region there is a simple connection between Green's function and the Riemann mapping function. For the unit disk $|w| < 1$ the Green's function with respect to the origin is evidently $-\log |w|$. Therefore, if $w = f(z)$ maps $\Omega$ onto the unit disk with $z_0$ going into the origin, we find by the invariance that
\begin{equation*}
g (z, z _ {0}) = - \log | f (z) |.
\end{equation*}
Conversely, if $g(z,z_{0})$ is known, the mapping function can be determined. The Green's function has an important symmetry property. Given two points $z_{1}, z_{2} \in \Omega$, we write for short $g(z,z_{1}) = g_{1}, g(z,z_{2}) = g_{2}$. By Theorem~\hyperref[thm:4.21]{21}, Chap.~\hyperref[chap:4]{4}, the differential $g_{1} * dg_{2} - g_{2} * dg_{1}$ is locally exact in the region obtained by omitting the points $z_{1}$ and $z_{2}$ from $\Omega$. If $c_{1}$ and $c_{2}$ are small circles about $z_{1}$ and $z_{2}$, described in the positive sense, the cycle $C - c_{1} - c_{2}$ is homologous to zero (as before, $C = C_{1} + \cdots + C_{n}$ ). Since $g_{1}$ and $g_{2}$ vanish on C, we conclude that
\begin{equation*}
\int_ {c _ {1} + c _ {2}} g _ {1} ^ {*} d g _ {2} - g _ {2} ^ {*} d g _ {1} = 0.
\end{equation*}
Introducing $G_{1} = g_{1} + \log |z - z_{1}|$ we have $^{*}dg_{1} = ^{*}dG_{1} - d\arg (z - z_{1})$ and find
\begin{equation*}
\begin{array}{r l} \int_ {c _ {1}} g _ {1} ^ {*} d g _ {2} - g _ {2} ^ {*} d g _ {1} = & \int_ {c _ {1}} G _ {1} ^ {*} d g _ {2} - g _ {2} ^ {*} d G _ {1} - \int_ {c _ {1}} \log | z - z _ {1} | ^ {*} d g _ {2} \\ & + \int_ {c _ {1}} g _ {2} d \arg (z - z _ {1}). \end{array}
\end{equation*}
On the right-hand side the first integral vanishes because $G_{1}$ and $g_{2}$ are harmonic inside $c_{1}$, and the second integral vanishes because $|z - z_{1}|$ is constant on $c_{1}$ and $*dg_{2}$ is an exact differential in a neighborhood of $z_{1}$. The last integral equals $2\pi g_{2}(z_{1})$ by the mean-value property of harmonic functions. In a symmetric way the integral over $c_{2}$ must equal $-2\pi g_{1}(z_{2})$, and it is proved that $g_{2}(z_{1}) - g_{1}(z_{2}) = 0$ or
\begin{equation*}
g (z _ {1}, z _ {2}) = g (z _ {2}, z _ {1}).
\end{equation*}
Because of this symmetry property the Green's function $g(z, z_0)$ is harmonic also in the second variable.


% Original PDF physical page 275
\sourcepage{275}
The conjugate function of $g(z,z_{0})$, denoted by $h(z,z_{0})$, is of course multiple-valued. It has above all the period $2\pi$ along a small circle c about $z_{0}$. In addition, it has the periods
\begin{equation*}
P _ {k} (z _ {0}) = \int_ {C _ {k}} d h (z, z _ {0}) = \int_ {C _ {k}} {}^{*}d g (z, z _ {0}) \quad (k = 1, \dots , n).
\end{equation*}
\begin{originalstatement}{Lemma}{3}\label{lem:6.3}
The period $P_{k}(z_{0})$ equals the harmonic measure $\omega_{k}(z_{0})$ multiplied by $2\pi$.
\end{originalstatement}
The proof is another application of Theorem~\hyperref[thm:4.21]{21}, Chap.~\hyperref[chap:4]{4}. We express the fact that the integral of $\omega_{k} * dg - g * d\omega_{k}$ over $C - c$ must vanish. The integral over $C$ reduces to $P_{k}(z_{0})$, and by the same computation as above the integral over $c$ equals $2\pi \omega_{k}(z_{0})$. Hence $P_{k}(z_{0}) = 2\pi \omega_{k}(z_{0})$.

\setcounter{section}{5}\setcounter{subsection}{2}\subsection{Parallel Slit Regions}\label{sec:6.5.3}
A little more explicitly than before, let us write
\begin{equation*}
g (z, z _ {0}) = G (z, z _ {0}) - \log | z - z _ {0} |\tag{16}\label{eq:6.16}
\end{equation*}
with $z_0 = x_0 + iy_0 \in \Omega$. We know that $G(z, z_0)$ is symmetric, and harmonic in each variable; as a function of $z$ it has the boundary values $\log |\xi - z_0|$.

Consider the difference quotient $Q(z,h) = (G(z,z_{0} + h) - G(z,z_{0}))/h$ where we choose h real and so small that $z_{0} + h$ is still in $\Omega$. This is a harmonic function of z with boundary values $(\log |\zeta - z_{0} - h| - \log |\zeta - z_{0}|)/h$. As $h \to 0$ these boundary values tend uniformly to $\partial/\partial x_{0} \log |\zeta - z_{0}| = -\operatorname{Re} 1/(\zeta - z_{0})$. It follows by the maximum-minimum principle that $Q(z,h)$ tends to its limit $(\partial/\partial x_{0})G(z,z_{0})$ uniformly, not only on compact sets, but on all of $\Omega$. If we include the boundary values, we have thus uniform convergence on the closure $\Omega^{-}$, which is a compact set. The conclusion is that $(\partial/\partial x_{0})G(z,z_{0})$ is harmonic in $\Omega$, as a function of z, and that it has the boundary values $-\operatorname{Re} 1/(\zeta - z_{0})$. If we compare with \hyperref[eq:6.16]{(16)} it follows that $u_{1}(z) = (\partial/\partial x_{0})g(z,z_{0})$ is harmonic for $z \neq z_{0}$, continuously zero on the boundary, and differs from $\operatorname{Re} 1/(z - z_{0})$ by a harmonic function.

The conjugate differential of $u_{1}(z)$ has certain periods $A_{k}$ along the contours $C_{k}$. But it is easy to construct a linear combination of $u_{1}(z)$ and the harmonic measures $\omega_{j}(z)$ whose conjugate differential is free from periods. Indeed, $u_{1} + \lambda_{1}\omega_{1} + \cdots + \lambda_{n-1}\omega_{n-1}$ has this property provided that
\begin{equation*}
\lambda_ {1} \alpha_ {1 k} + \lambda_ {2} \alpha_ {2 k} + \dots \lambda_ {n - 1} \alpha_ {n - 1, k} = - A _ {k} \quad (k = 1, \dots , n - 1).
\end{equation*}
We know already that this inhomogeneous system of equations always has a solution. We have thus established the existence of a function
% Original PDF physical page 276
\index{Function!Green's|)}\index{Green's function|)}\sourcepage{276}
$p(z)$\index{Mapping!slit|(} which is single-valued and analytic in $\Omega$, except for a simple pole with the residue 1 at $z_{0}$, and whose real part is constant on each contour. By these requirements $p(z)$ is uniquely determined up to an additive constant.

By differentiation with respect to $y_{0}$ we conclude quite similarly that $v_{2}(z) = -(\partial/\partial y_{0})g(z,z_{0})$ vanishes on the boundary and has the same singularity as $\operatorname{Im} 1/(z - z_{0})$. If a suitable linear combination of harmonic measures is added, the conjugate function becomes single-valued. Hence there exists a single-valued analytic function $q(z)$ with the singular part $1/(z - z_{0})$ whose imaginary part is constant on each contour.

The functions $p(z)$ and $q(z)$ lead to simple canonical mappings\index{Canonical mapping}.

\begin{originalstatement}{Theorem}{11}\label{thm:6.11}
The mappings determined by $p(z)$ and $q(z)$ are one to one, and the image of $\Omega$ is a slit region whose complement consists of $n$ vertical or horizontal segments, respectively (Fig. 6-5a, b).
\end{originalstatement}
The proof is quite similar to that of Theorem~\hyperref[thm:6.10]{10}. This time the expression
\begin{equation*}
\sum_ {k = 1} ^ {n} \frac {1}{2 \pi i} \int_ {C _ {k}} \frac {p ^ {\prime} (z) d z}{p (z) - w _ {0}}\tag{17}\label{eq:6.17}
\end{equation*}
represents the number of zeros of $p(z) - w_{0}$ minus the number of poles. But it is easy to see that \hyperref[eq:6.17]{(17)} vanishes for all $w_{0}$, including boundary values. In the latter case the principal value must be formed, but if $w_{0}$ is taken on $C_{k}$ the imaginary part of $p' dz/(p - w_{0})$ vanishes along $C_{k}$ and there is no difficulty whatsoever. Since there is exactly one pole we conclude that $p(z)$ takes every value once in the interior of $\Omega$, twice on the boundary, or once on the boundary with the multiplicity 2. The rest of the proof is an exact duplication of the earlier reasoning. The proof remains valid for $q(z)$ without change.

Parallel slit regions may be thought of as canonical regions, but they are not all conformally inequivalent, even if it is required that the point at $\infty$ should correspond to itself. For instance, the mappings by $p(z)$

\bookfigure{276}{1}{FIG. 6-5. Parallel slit regions.}

% Original PDF physical page 277
\index{Mapping!slit|)}\sourcepage{277}
and $iq(z)$ lead to vertical slit regions which are different, but conformally equivalent. It is only for mappings with the same residue at $z_{0}$ that the slit mappings are uniquely determined, except for a parallel translation.

\subsection*{EXERCISES}
\originalnumber{1}{exercise:6.5.3.8.1}
1. Prove that $g(z,z_{0})$ is simultaneously continuous in both variables, for $z \neq z_{0}$. Hint: Apply the maximum-minimum principle to $G(z,z_{0})$.

\originalnumber{2}{exercise:6.5.3.8.2}
2. Show that the function $e^{-i\alpha}(q \cos \alpha + ip \sin \alpha)$ maps $\Omega$ onto a region bounded by inclined slits.

\originalnumber{3}{exercise:6.5.3.8.3}
*3. Using Ex.~\hyperref[exercise:6.5.3.8.2]{2}, show that $p + q$ maps $\Omega$ in a one-to-one manner onto a region bounded by convex contours. Comments:

(i) A closed curve is said to be convex if it intersects every straight line at most twice.

(ii) To prove that the image of $C_{k}$ under $p + q$ is convex we need only show that for every $\alpha$ the function $\operatorname{Re}(p + q)e^{i\alpha}$ takes no value more than twice on $C_{k}$. But $\operatorname{Re}(p + q)e^{i\alpha}$ differs from $\operatorname{Re}(q\cos\alpha + ip\sin\alpha)$ only by a constant, and the desired conclusion follows by the properties of the mapping function in Ex.~\hyperref[exercise:6.5.3.8.2]{2}.

(iii) Finally, the argument principle can be used to show that the images of the contours $C_{k}$ have winding number 0 with respect to all values of $p + q$. This implies, in particular, that the convex curves lie outside of each other.


% Original PDF physical page 278
\index{Canonical mapping|)}\index{Conformal mapping|)}\index{Mapping!conformal|)}\sourcepage{278}
